\begin{theorem}
For each $\iota>0$ and $k\in\mathbb Z^+$, there is $D_0$ such that if
$D\ge D_0$ and $\mathcal H$ is a $k$-uniform $k$-simple hypergraph with
maximum degree at most $D$, then
\[
\chi'(\mathcal H)\le
\left(1-\frac{k-1}{2k^2}
+\frac{(k-1)(k-2)}{6k^3\sqrt{k}}+\iota\right)kD.
\]
\end{theorem}
\begin{proof}
Fix $\iota>0$ and a positive integer $k$.  First consider the case
$k=1$.  Then every edge of $\mathcal H$ consists of a single vertex by
\noderef{UniformHypergraph}.  For each vertex $v$, at most $D$ edge indices
are incident with $v$ by \noderef{MaxDegreeAtMost} and
\noderef{HypergraphDegree}.  Color the edges incident with each fixed
vertex injectively from a common set of $D$ colors, reusing the same color
set at different vertices.  Edges supported on different vertices are
disjoint, and edges supported on the same vertex receive distinct colors, so
this is a proper edge-coloring in the sense of
\noderef{ProperEdgeColoring}.  Hence
\noderef{ChromaticIndexAtMostReal} gives $\chi'(\mathcal H)\le D$.  When
$k=1$, the coefficient in the theorem is $1+\iota$, so
$D\le(1+\iota)D$, and the desired bound follows.

We may therefore assume $k\ge2$.  Put
\[
c_k=
1-\frac{k-1}{2k^2}
+\frac{(k-1)(k-2)}{6k^3\sqrt{k}}.
\]
For $k\ge2$ the elementary inequality $c_k<1$ holds: the positive term is
strictly smaller than $(k-1)/(2k^2)$ after clearing denominators and using
$k-2<3k\sqrt{k}$ for $k\ge2$. Also $c_k>0$: the subtracted term
$(k-1)/(2k^2)$ is less than $1$, and the final term is nonnegative.
Set $\varepsilon=\min\{\iota/2,(1-c_k)/2\}>0$ and
$\theta_*=c_k+\varepsilon\in(0,1)$. Then $2\varepsilon\le\iota$.
Let $D_g$ be the natural threshold from \noderef{GeneralColoringTheorem}
for $\varepsilon,k$, and let $D_b$ be the threshold from
\noderef{KUniformKSimpleBParameterEstimate} for error $\varepsilon$ and $k$.
Choose a natural number $N>3(D_b+1)/\varepsilon$ and put
$D_0=\max\{D_g,N\}$. In particular, for every natural $D\ge D_0$,
$\varepsilon D/3>D_b+1$.

Now let $D\ge D_0$ and let $\mathcal H\in H(k,k,D)$.  By
\noderef{HypergraphClass}, this means that $\mathcal H$ is $k$-uniform,
$k$-simple, and has maximum degree at most $D$; these are precisely the
three components \noderef{UniformHypergraph}, \noderef{TSimpleHypergraph},
and \noderef{MaxDegreeAtMost}.  We verify the hypothesis needed to apply
\noderef{GeneralColoringTheorem} to $\mathcal H$ with parameter
$\theta_*$.

Let $D'=r\ge \varepsilon D/3$ be real, and let
$H\subseteq\mathcal H$ be a sub-hypergraph with
$\Delta(H)\le D'$.  The relation $H\subseteq\mathcal H$ is
\noderef{SubhypergraphOf}: there is an injective map from the edge indices
of $H$ to those of $\mathcal H$ preserving their vertex sets.
By \noderef{SubhypergraphInheritance}, $H$ is still $k$-uniform and
$k$-simple. In particular, distinct child edge indices have distinct
ambient representatives, so the ambient simplicity bound applies.
Set $n=\lfloor r\rfloor_+$. By \noderef{RealDegreeFloorBound}, the real
degree bound implies \noderef{MaxDegreeAtMost} at the natural bound $n$.
Also $r\ge D_b+1\ge1$, so $n\ge D_b$ and $n\ge1$.
Thus \noderef{HypergraphClass} places $H$ in $H(k,k,n)$, and
\noderef{KUniformKSimpleBParameterEstimate} applies at $n$.
For every edge
$e\in E(H)$,
\[
b^{\mathbb R}_{L(H),kr}(e)
\le b_{L(H),kn}(e)
\le
1-\frac{k-1}{2k^2}
+\frac{(k-1)(k-2)}{6k^3\sqrt{k}}+\varepsilon
=\theta_* .
\]
The first inequality is \noderef{LocalBParameterRealFloorTransfer};
its hypotheses hold because $k>0$, $H$ is $k$-uniform, $r\ge1$, and
all vertex degrees are at most $r$.
This is exactly the local $b$-parameter hypothesis in
\noderef{GeneralColoringTheorem}; the graph in that hypothesis is
\noderef{LineGraphOfHypergraph} and the parameter is
\noderef{LocalBParameterReal}. We apply the real main parameter of that
theorem to the real cast of our original natural $D$.

Applying \noderef{GeneralColoringTheorem} gives
\[
\chi'(\mathcal H)\le(\theta_*+\varepsilon)kD.
\]
Finally $2\varepsilon\le\iota$, so
\[
\theta_*+\varepsilon
=
1-\frac{k-1}{2k^2}
+\frac{(k-1)(k-2)}{6k^3\sqrt{k}}+2\varepsilon
\le
1-\frac{k-1}{2k^2}
+\frac{(k-1)(k-2)}{6k^3\sqrt{k}}+\iota.
\]
Substituting this inequality into the preceding coloring bound gives the
claimed estimate.
\end{proof}
